\begin{proposition}[5.12.5]
\label{IV.2.s5.12.5-fr}
Sean $A$ un anillo local noetheriano catenario, $t$ un elemento regular de $A$
perteneciente al ideal maximal de $A$, y $k$ un entero $\geq0$. Si el anillo
$A/tA$ es reducido, equidimensional y verifica la propiedad $(R_k)$, el anillo
$A$ posee también estas tres propiedades.
\end{proposition}

Ya se sabe por (3.4.6) que $A$ es reducido, y de (5.12.2), aplicado para
$k=1$, se deduce que $A$ es equidimensional. Pongamos
$X=\operatorname{Spec}(A)$, $X_0=V(t)=\operatorname{Spec}(A/tA)$, y sea $Z$
el conjunto de los puntos $x\in X$ en los que $X$ no es regular; en virtud de
(\textbf{0}, 17.3.2), toda especialización de un punto de $Z$ pertenece a $Z$. Por otra
parte, de (\textbf 0, 17.1.8) se deduce que, para todo punto $x\in X_0$ en el que $X_0$
es regular, $X$ también es regular; por tanto, el conjunto $Z'_0$ de los puntos
en los que $X_0$ no es regular contiene $Z_0=Z\cap X_0$; la hipótesis implica,
pues, que
\[
k\leq\operatorname{codim}(Z'_0,X_0)\leq\operatorname{codim}(Z_0,X_0).
\]
Pero, puesto que $A$ es equidimensional y catenario, se tiene, por (5.12.1),
\[
\operatorname{codim}(Z_0,X_0)\leq\operatorname{codim}(Z,X)
\]
lo que prueba que $X$ verifica $(R_k)$.

\begin{remark}[5.12.6]
\label{IV.2.s5.12.6-fr}
Si, en el enunciado de (5.12.5), se supone únicamente que $A/tA$ sea reducido y
posea la propiedad $(R_k)$, no se sigue necesariamente que $A$ posea la
propiedad $(R_k)$. Por ejemplo, sean $k$ un cuerpo, $P(U,V)$ un polinomio
irreducible de $k[U,V]$ tal que la curva $\Gamma$ definida por el ideal
principal $(P)$ en el plano afín $\operatorname{Spec}(k[U,V])$ tenga un punto
singular correspondiente al ideal maximal $(U)+(V)$ (por ejemplo
$P(U,V)=U(U^2+V^2)+(U^2\setmin V^2)$, que da una cúbica con punto doble). Sean entonces
$B$ el anillo de polinomios en cuatro indeterminadas $k[T,U,V,W]$, $C$ el anillo
local de $B$ en el ideal maximal $BT+BU+BV+BW$, y sean
\[
\mathfrak p=CW+CP(U\setmin T,V),\qquad \mathfrak q=CU.
\]
Consideremos entonces el anillo local $A=C/\mathfrak p\mathfrak q$
(geométricamente, si $X$ es el subpreesquema cerrado del espacio afín de cuatro
dimensiones, formado por un hiperplano $H$ y un «cilindro» $L$ de dos dimensiones
que tiene por «base» la curva $\Gamma$ y cuya «recta singular» $Y$ no está
contenida en el hiperplano $H$, entonces $A$ es el anillo local de $X$ en el
punto $x$ donde $Y$ corta a $H$). Se ve entonces inmediatamente que
$\operatorname{Spec}(A/tA)$ (donde $t$ es la imagen canónica de $T$ en $A$)
tiene como único punto singular a $x$, cuyo anillo local es el propio $A/tA$,
que es reducido y de dimensión 2; en cambio, el punto genérico $y$ de la
«recta singular» $Y$ (definida por la imagen en $A$ del ideal
$C(U\setmin T)+CV+CW$) es un punto singular de $X$, y
$\mathcal O_{X,y}$ es de dimensión 1; dicho de otro modo, $A/tA$ es reducido y
verifica $(R_1)$, pero $A$ no verifica $(R_1)$.
\end{remark}

\oldpage[IV-2]{129}

\begin{corollary}[5.12.7]
\label{IV.2.s5.12.7-fr}
Sean $A$ un anillo local noetheriano catenario, $t$ un elemento regular de
$A$ perteneciente al ideal maximal de $A$. Si $A/tA$ es íntegro e
integralmente cerrado, lo mismo ocurre con $A$.
\end{corollary}

En virtud del criterio de Serre (5.8.6) y del hecho de que el anillo $A$ es
local, basta mostrar que $\operatorname{Spec}(A)$ verifica $(S_2)$ y $(R_1)$.
Pero la hipótesis sobre $A/tA$ implica que $A$ verifica $(R_1)$ por (5.12.5), y
que $A$ verifica $(S_2)$ por (5.12.4).
